\subsection{Đề 2}
\Opensolutionfile{ans}[ans/ans-2-Otc5-De2]
\caulc
\begin{ex}%[2H5N2-2]
	Trong không gian $Oxyz$, cho đường thẳng $\heva{&x=1-t\\&y=-2+2t\\&1+t}$. Véc-tơ nào dưới đây là vectơ chỉ phương của $d$?
	\choice
	{$(1;-2;1)$}
	{$(1;2;1)$}
	{$(-1;-2;1)$}
	{\True $(-1;2;1)$}
	\loigiai{
		Dựa vào phương trình tham số của đường thẳng $d$ ta có véc-tơ chỉ phương của $d$ là $\overrightarrow{n}=(-1;2;1)$.
	}
\end{ex}

\begin{ex}%[2H5H1-5]
	Trong không gian $Oxyz$, cho $(P) \colon 3x-4y+2z+4=0$ và điểm $A(1;-2;3)$. Tính khoảng cách từ $A$ đến $(P)$.
	\choice
	{\True $\dfrac{21}{\sqrt{29}}$}
	{$\dfrac{\sqrt{5}}{3}$}
	{$\dfrac{5}{\sqrt{29}}$}
	{$\dfrac{5}{9}$}
	\loigiai{
		$\mathrm{d}(A,(P))=\dfrac{ \left| 3 \cdot 1 - 4 \cdot (-2) + 2 \cdot 3 + 4 \right|}{\sqrt{3^2+(-4)^2+2^2}} = \dfrac{21}{\sqrt{29}}.$
	}
\end{ex}

\begin{ex}%[2H5N3-2]
	Trong không gian $Oxyz$, tìm tâm và bán kính của mặt cầu có phương trình $(x-1)^2+(y+4)^2+(z-3)^2=18$.
	\choice
	{$I(-1;-4;3)$, $R=\sqrt{18}$}
	{$I(1;-4;-3)$, $R=\sqrt{18}$}
	{$I(1;4;3)$, $R=\sqrt{18}$}
	{\True $I(1;-4;3)$, $R=\sqrt{18}$}
	\loigiai{
		Mặt cầu có tâm $I(1;-4;3)$, $R=\sqrt{18}$.
	}
\end{ex}

\begin{ex}%[2H2N2-3]
	Trong không gian $Oxyz$, cho $\overrightarrow{OA}=\overrightarrow{i}-2\overrightarrow{j}+3\overrightarrow{k}$. Tìm tọa độ điểm $A$.
	\choice
	{$A\left(-1;-2;-3\right)$}
	{$A\left(1;2;3\right)$}
	{$A\left(2;-4;6\right)$}
	{\True $A\left(1;-2;3\right)$}
	\loigiai{
		\begin{center}
			$\overrightarrow{OA}=\overrightarrow{i}-2\overrightarrow{j}+3\overrightarrow{k}=(1;-2;3)\Rightarrow A(1;-2;3)$.
		\end{center}
	}
\end{ex}
\begin{ex}%[2H2H1-6]
	Trong không gian $Oxyz$, cho hai mặt phẳng $ (P)\colon -6x+my-2mz-m^{2}=0 $ và $ (Q)\colon 2x+y-2z+3=0 $ ($ m $ là tham số). Tìm $ m $ để mặt phẳng $ (P) $ vuông góc với mặt phẳng $ (Q) $.
	\choice
	{$ m=12 $}
	{$ m=\dfrac{12}{7} $}
	{$ m=\dfrac{5}{12} $}
	{\True $ m=\dfrac{12}{5} $}
	\loigiai{
		Véc-tơ pháp tuyến của $ (P) $ là $ \overrightarrow{n}_{P}=(-6;m;-2m) $, véc-tơ pháp tuyến của $ (Q) $ là $ \overrightarrow{n}_{Q}=(2;1;-2) $.\\
		Mặt phẳng $ (P) $ vuông góc với mặt phẳng $ (Q) $ khi và chỉ khi
		\[
		\overrightarrow{n}_{P}\cdot\overrightarrow{n}_{Q}=0\Leftrightarrow -12+m+4m=0\Leftrightarrow m=\dfrac{12}{5}.
		\]
	}
\end{ex}

\begin{ex}%[2H2H2-3]
	Trong không gian $Oxyz$, cho hai véc-tơ $\overrightarrow{a}=(-4;5-3)$, $\overrightarrow{b}=(2;-2;1)$. Tìm tọa độ của véc-tơ $\overrightarrow{x}=\overrightarrow{a}+2\overrightarrow{b}$.
	\choice
	{$\overrightarrow{x}=(0;-1;1)$}
	{$\overrightarrow{x}=(-8;9;1)$}
	{$\overrightarrow{x}=(2;3;-2)$}
	{\True $\overrightarrow{x}=(0;1;-1)$}
	\loigiai{
		Ta có $\overrightarrow{x}=\overrightarrow{a}+2\overrightarrow{b}=(0;1;-1)$.}
\end{ex}

\begin{ex}%[2H2H2-3]
	Trong không gian $Oxyz$, cho tam giác $ABC$ có $A(1;3;5),~B(2;0;1)$ và $G(1;4;2)$ là trọng tâm. Tìm tọa độ điểm $C$.
	\choice
	{$C\left(\dfrac{4}{3};\dfrac{7}{3};\dfrac{8}{3}\right)$}
	{$C(0;0;9)$}
	{\True $C(0;9;0)$}
	{$C(0;-9;0)$}
	\loigiai
	{
		Do $G(1;4;2)$ là tọa độ trọng tâm tam giác $ABC$, ta có $\heva{&x_C=3x_G-(x_A+x_B)=0\\&y_C=3y_G-(y_A+y_B)=9\\&z_C=3z_G-(z_A+z_B)=0.}$\\
		Vậy $C(0;9;0)$.
	}
\end{ex}

\begin{ex}%[2H3B2-4]
	Trong không gian $Oxyz$, cho điểm $A(3;-1;1)$. Hình chiếu vuông góc của $A$ trên mặt phẳng $(Oxy)$ là điểm
	\choice
	{\True $N(3;-1;0)$}
	{$M(3;0;0)$}
	{$P(0;-1;0)$}
	{$Q(0;0;1)$}
	\loigiai{
		Hình chiếu của $A(3;-1;1)$ trên mặt phẳng $(Oxy)$ có tọa độ là $(3;-1;0)$.
	}
\end{ex}
\begin{ex}%[2H5V1-3]
	Trong không gian $Oxyz$, mặt phẳng $(\alpha)$ đi qua hai điểm $A(1;2;-2), B(2;-1;4)$ và vuông góc với mặt phẳng $(\beta)\colon x-2y-z+1=0$ có phương trình là
	\choice
	{\True $15x+7y+z-27=0$}
	{$15x-7y+z-27=0$}
	{$15x-7y+x+27=0$}
	{$15x+7y-z+27=0$}
	\loigiai{
		Ta có $\overrightarrow{AB}=(1;-3;6),\overrightarrow{n}_\beta=(1;-2;-1)$.\\
		Gọi $\overrightarrow{n}_\alpha$ là một véc-tơ pháp tuyến của $(\alpha)$.\\
		Vì $(\alpha)$ đi qua $A,B$ và vuông góc với $(\beta)$ nên $\overrightarrow{n}_\alpha\perp\overrightarrow{AB}$ và $\overrightarrow{n}_\alpha\perp\overrightarrow{n}_\beta$ mà $\overrightarrow{AB}$ không cùng phương với $\overrightarrow{n}_\beta$, do đó $$\overrightarrow{n}_\beta=\left[\overrightarrow{AB},\overrightarrow{n}_\beta\right]=(15;7;1).$$
		Vậy, phương trình mặt phẳng $(\alpha)$ đi qua $A(1;2;-2)$ và có một véc-tơ pháp tuyến $\overrightarrow{n}_\alpha=(15;7;1)$ là $$15(x-1)+7(y-2)+1(z+2)=0\Leftrightarrow 15x+7y+z-27=0.$$
	}
\end{ex}

\begin{ex}%[2H5V2-3]
	Trong không gian $Oxyz$, phương trình mặt phẳng cách đều hai đường thẳng $d_1 \colon \dfrac{x+1}{2}=\dfrac{y-1}{3}=\dfrac{z-2}{1}$ và $d_2 \colon \dfrac{x-2}{1}=\dfrac{y+2}{5}=\dfrac{z}{-2}$ là
	\choice
	{$-11x+5y+7z+11=0$}
	{\True $-11x+5y+7z+1=0$}
	{$11x-5y-7z+1=0$}
	{$-11x+5y+7z-1=0$}
	\loigiai{
		Đường thẳng $d_1$ đi qua điểm $A(-1;1;2)$ và có véc-tơ chỉ phương là $\overrightarrow{u_1}=(2;3;1)$.\\
		Đường thẳng $d_2$ đi qua điểm $B(2;-2;0)$ và có véc-tơ chỉ phương là $\overrightarrow{u_2}=(1;5;-2)$.\\
		Ta có $\left[ \overrightarrow{u}_1,\overrightarrow{u}_2\right] =(-11;5,7)$. \\
		Mặt phẳng $\alpha$ cách đều hai đường thẳng $d_1$, $d_2$ có dạng $-11x+5y+7z+m=0$.\\
		Ta có $\mathrm{d}(A,\alpha)=\mathrm{d}(B,\alpha) \Leftrightarrow \dfrac{|30+m|}{\sqrt{11^2+5^2+7^2}}=\dfrac{|-32+m|}{\sqrt{11^2+5^2+7^2}} \Leftrightarrow m=1$.\\
		Vậy mặt phẳng đó có phương trình là $-11x+5y+7z+1=0$.}
\end{ex}

\begin{ex}%[2H5C3-3]
	Trong không gian $Oxyz$, cho mặt cầu $(S)\colon x^2+y^2+z^2-2x-2y+4z-1=0$ và mặt phẳng $(P)\colon x+y-z-m=0$. Tìm tất cả $m$ để $(P)$ cắt $(S)$ theo giao tuyến là một đường tròn có bán kính lớn nhất.
	\choice
	{$m=7$}
	{\True $m=4$}
	{$m=0$}
	{$m=-4$}
	\loigiai
	{
		Mặt cầu $(S)$ có tâm $I(1;1;-2)$. Để $(P)$ cắt $(S)$ theo đường tròn có bán kính lớn nhất thì đó là đường tròn lớn. Suy ra $I \in (P) \Leftrightarrow m=4$.
	}
\end{ex}

\begin{ex}%[2H5V3-]
	Trong không gian $Oxyz$, mặt phẳng $(P)\colon 2x+6y+z-3=0$ cắt trục $Oz$ và đường thẳng $d\colon\dfrac{x-5}{1}=\dfrac{y}{2}=\dfrac{z-6}{-1}$ lần lượt tại $A$ và $B$. Phương trình mặt cầu đường kính $AB$ là
	\choice
	{$(x+2)^2+(y-1)^2+(z+5)^2=9$}
	{$(x-2)^2+(y+1)^2+(z-5)^2=36$}
	{\True $(x-2)^2+(y+1)^2+(z-5)^2=9$}
	{$(x+2)^2+(y-1)^2+(z+5)^2=36$}
	\loigiai{
		\begin{itemize}
			\item Mặt phẳng $(P)\colon 2x+6y+z-3=0$ cắt trục $Oz$ tại điểm $A(0;0;3)$.
			\item Tọa độ $B=d\cap (P)$ thỏa mãn hệ
			\[\heva{&\dfrac{x-5}{1}=\dfrac{y}{2}=\dfrac{z-6}{-1}=t\\&2x+6y+z-3=0}
			\Leftrightarrow\heva{&x=t+5\\&y=2t\\&z=-t+6\\&2(t+5)+6\cdot 2t-t+6-3=0}\Leftrightarrow\heva{&t=-1\\&x=4\\&y=-2\\&z=7}\Rightarrow B(4;-2;7).\]
			\item Tâm mặt cầu là trung điểm $I$ của đoạn $AB$, suy ra $I(2;-1;5)$. Bán kính mặt cầu là $R=\dfrac{1}{2}AB=3$.
			\item Vậy phương trình mặt cầu là $(x-2)^2+(y+1)^2+(z-5)^2=9$.
		\end{itemize}
	}
\end{ex}

\Closesolutionfile{ans}
\indapan{6}{ans/ans-2-Otc5-De2}

\cauds
\Opensolutionfile{ans}[ans/ans-2-Otc5-De2-DS]

\begin{ex}%[2H5V1-5]
Trong không gian $Oxyz$, cho điểm $A(2;1;0)$ và mặt phẳng $(P)\colon 2x+y-2z+1=0$.
	\choiceTF
	{\True Một véc-tơ pháp tuyến của mặt phẳng $(P)$ là $\vec{n}=(2;1;-2)$}
	{Điểm $A$ thuộc mặt phẳng $(P)$}
	{Khoảng cách từ điểm $A$ đến mặt phẳng $(P)$ bằng 3}
	{\True Phương trình mặt cầu tâm $A$ và tiếp xúc với mặt phẳng $(P)$ là $(x-2)^2+(y-1)^2+z^2=9$}

	\loigiai{
		\begin{itemchoice}
			\itemch Đúng. Vì $\vec{n}=(2;1;-2)$
			\itemch Sai. Vì $2\cdot 2+1-2\cdot 0+1=6\neq 0$ nên điểm $A$ không thuộc $(P)$.
			\itemch Sai. Vì $d\left(A;(P)\right)=2$.
			\itemch Đúng. Vì $R=d(A;(P))=2$ nên $(S):(x-2)^2+(y-1)^2+(z+2)^2=4$.
		\end{itemchoice}
	}
\end{ex}
\begin{ex}%[2H5V2-6]
	Trong không gian $Oxyz$, cho mặt phẳng $(P)\colon x-y+z-1=0$ và mặt cầu $(S)\colon (x-1)^2+y^2+(z+1)^2=25$.
	\choiceTF
	{\True Một véc-tơ pháp tuyến của mặt phẳng $(P)$ là $\vec{n}=\left(1;-1;1\right)$}
	{\True Mặt cầu $(S)$ có tâm $I\left(1;0;-1\right)$}
	{\True Mặt phẳng $(P)$ cắt mặt cầu $(S)$ theo giao tuyến là hình tròn}
	{Bán kính của hình tròn giao tuyến tạo bởi $(P)$ và $(S)$ bằng 5}
	\loigiai{
		\begin{itemchoice}
			\itemch Đúng. Vì $\vec{n}=(1;-1;1)$.
			\itemch Đúng. Vì $I(1;0;-1)$
			\itemch Đúng. Vì $d\left(I;(P)\right)=\dfrac{\sqrt{3}}{3}<R$.
			\itemch Sai. Vì $r^2=R^2-d^2=25-\dfrac{1}{3}=\dfrac{74}{3}$
		\end{itemchoice}
	}
\end{ex}
\begin{ex}%[2H5V2-6]
	Trong không gian $Oxyz$, cho điểm $A(1;-1;2)$ và đường thẳng $d:\dfrac{x-1}{2}=\dfrac{y+2}{1}=\dfrac{z-4}{-3}$.
	\choiceTF
	{Véc-tơ $\vec{u}=(2;1;3)$ là một véc-tơ chỉ phương của đường thẳng $d$}
	{\True Điểm $A$ không thuộc đường thẳng $d$}
	{Khoảng cách từ điểm $A$ đến đường thẳng $d$ bằng 6}
	{\True Phương trình đường thẳng đi qua $A$ và song song với đường thẳng $d$ là $\dfrac{x-1}{2}=\dfrac{y+1}{1}=\dfrac{z-2}{-3}$}
	
	\loigiai{
		\begin{itemchoice}
			\itemch Sai. Vì $\vec{u}=(2;1;-3)$.
			\itemch Đúng. Vì $\dfrac{1-1}{2}=\dfrac{-1+2}{1}=\dfrac{2-4}{-3}$ vô lý.
			\itemch Sai. Vì $d(A;d)=\dfrac{\sqrt{2}}{2}$.
			\itemch Sai. Vì  $d:\dfrac{x-1}{2}=\dfrac{y+1}{1}=\dfrac{z-2}{-3}$.
		\end{itemchoice}
	}
\end{ex}
\begin{ex}%[2H5V2-7]
	Trong không gian $Oxyz$, cho đường thẳng $d\colon\heva{&x=2+t\\&y=-4+3t\\&z=1+t}$ và mặt phẳng $(P)\colon3x+y-z+1=0$.
	\choiceTF
	{\True Một véc-tơ chỉ phương của đường thẳng $d$ là $\vec{u}=(1;3;1)$}
	{\True Một véc-tơ pháp tuyến của mặt phẳng $(P)$ là $\vec{n}=(3;1;-1)$}
	{Giá trị $\vec{u_d}\cdot\vec{n_{(P)}}=12$}
	{Góc giữa đường thẳng $d$ và mặt phẳng $(P)$ bằng $45^o$}
	\loigiai{
		\begin{itemchoice}
			\itemch Đúng. Vì $\vec{u_d}=(1;3;1).$
			\itemch Đúng. Vì $\vec{n_{(P)}}=(3;1;-1)$.
			\itemch Sai. Vì $\vec{u_d}.\vec{n_{(P)}}=3+3-1=8$.
			\itemch Sai. Vì $\sin \left(d;(P)\right)=\dfrac{8}{11}$.
		\end{itemchoice}
	}
\end{ex}


\Closesolutionfile{ans}
\indapan{3}{ans/ans-2-Otc5-De2-DS}

\Opensolutionfile{ans}[ans/ans-2-Otc5-De2-KQ]
\caukq
\begin{ex}%[2H2N1-3]
	Trong không gian $Oxyz$, cho hai véc-tơ $\overrightarrow{a}=(2;1;0)$, $\overrightarrow{b}=(-1;0;2)$. Tính $\cos(\overrightarrow{a},\overrightarrow{b})$.
\shortans{$0,4$}
	\loigiai{
		$$\cos(\overrightarrow{a},\overrightarrow{b})=\dfrac{\overrightarrow{a}\cdot\overrightarrow{b}}{|\overrightarrow{a}|\cdot|\overrightarrow{b}|}
		=\dfrac{2\cdot(-1)+1\cdot0+0\cdot 2}{\sqrt{2^2+1^2+0^2}\cdot\sqrt{(-1)^2+0^2+2^2}}=-\dfrac{2}{5}.$$
	}
\end{ex}
\begin{ex}%[2H5N2-5]
	Trong không gian $Oxyz$, cho đường thẳng $d\colon\heva{&x=(m-1)t\\&y=(2m+1)t\\&z=1+(2m^2+1)t}$. Với giá trị nào của $m$ thì đường thẳng $d$ nằm trong mặt phẳng $Oyz$?
\shortans{$1$}
	\loigiai{
		Đường thẳng $d$ có một véc-tơ chỉ phương là $\overrightarrow{u}_d=(m-1;2m+1;2m^2+1)$.\\
		Vì $d$ nằm trong mặt phẳng $(Oyz)$ nên $$\overrightarrow{u}_d\perp\overrightarrow{i}\Leftrightarrow\overrightarrow{u}_d\cdot\overrightarrow{i}=0\Leftrightarrow(m-1)\cdot1+(2m-1)\cdot0+(2m^2+1)\cdot0=0\Leftrightarrow m=1.$$
		Thử lại, $m=1\Rightarrow d\colon \heva{&x=0\\&y=3t\\&z=1+3t.}$\\
		Lấy $M(0;0;1)\in d$, dễ thấy $M\in (Oyz)\Rightarrow$ $d$ nằm trong mặt phẳng $(Oyz)$.\\
		Vậy $m=1$ là giá trị cần tìm.
	}
\end{ex}

\begin{ex}%[2H5N1-6]
	Trong không gian $Oxyz$, cho hai mặt phẳng $(P) \colon x + (m + 1)y - 2z + m = 0$ và $(Q) \colon 2x - y + 3 = 0$ với $m$ là tham số thực. Tìm $m$ để $(P)$ vuông góc với $(Q)$.
\shortans{$1$}
	\loigiai{
		$\bullet$ $(P)$ có vec-tơ pháp tuyến là $\overrightarrow{n_P} = (1; m + 1; -2)$ và $(Q)$ có vec-tơ pháp tuyến $\overrightarrow{n_Q} = (2; -1; 0)$. \\
		$\bullet$ $(P) \perp (Q) \Leftrightarrow \overrightarrow{n_P} \perp \overrightarrow{n_Q} \Leftrightarrow 1\cdot 2  + (m + 1) \cdot (-1) + (-2) \cdot 0 = 0 \Leftrightarrow m = 1$.
	}
\end{ex}
\begin{ex}%[2H5V2-6]
	Trong không gian $Oxyz$, cho hai điểm $A(0;-1;2),\ B(1;1;2)$ và đường thẳng $d\colon \dfrac{x+1}{1}=\dfrac{y}{1}=\dfrac{z-1}{1}$. Biết điểm $M(a;b;c)$ thuộc đường thẳng $d$ sao cho tam giác $MAB$ có diện tích nhỏ nhất. Khi đó, giá trị $T=a+2b+3c$ bằng 
	\shortans{$10$}
	\loigiai{
		Vì $S_{\text{MAB}}=\dfrac{1}{2}\cdot AB\cdot\mathrm{d}(M,AB)$ nên $S_{\text{MAB}}$ nhỏ nhất khi $\mathrm{d}(M,AB)$ nhỏ nhất.\\
		Phương trình của $AB\colon \heva{&x=t\\ &y=-1+2t\\ &z=2}$. Dễ dàng kiểm tra $AB$ và $d$ chéo nhau.\\
		Gọi $H$ là hình chiếu của $M$ lên đường thẳng $AB$. Khi đó $\mathrm{d}(M,AB)=MH$ nhỏ nhất khi $MH$ là đoạn vuông góc chung của $d$ và $AB$.\\
		Ta có\\
		$M\in d\Rightarrow M(-1+s;s;1+s)$, $H\in AB\Rightarrow H(t;-1+2t;2)$, $\overrightarrow{MH}=(t-s+1;2t-s-1;1-s)$.\\
		Véc-tơ chỉ phương của $d$ và $AB$ theo thứ tự là $\overrightarrow{u}=(1;1;1)$, $\overrightarrow{v}=(1;2;0)$.\\
		Vì $\heva{&\overrightarrow{MH}\perp \overrightarrow{u}\\ &\overrightarrow{MH}\perp \overrightarrow{v}}\Leftrightarrow \heva{& 1(t-s+1)+1(2t-s-1)+1(1-s)=0\\ & 1(t-s+1)+2(2t-s-1)+0(1-s)=0}\Leftrightarrow \heva{& t=1\\ & s=\dfrac{4}{3}}$.\\
		Vậy $M\left(\dfrac{1}{3};\dfrac{4}{3};\dfrac{7}{3}\right)$. Từ đó $T=\dfrac{1}{3}+2\cdot \dfrac{4}{3}+3\cdot \dfrac{7}{3}=10$.
	}
\end{ex}
\begin{ex}%[2H5V1-5]
	Trong không gian $Oxyz$, cho mặt phẳng $(P)\colon x-z+10=0$ và điểm $A(1; 0;0)$. Mặt phẳng $(\alpha)$ đi qua $A$, vuông góc với $(P)$, cách gốc tọa độ $O$ một khoảng bằng $\dfrac{2}{3}$ và cắt các tia $Oy$, $Oz$ lần lượt lại các điểm $B$, $C$ không trùng $O$. Thể tích khối tứ diện $OABC$ bằng (kết quả làm tròn đến hàng phần trăm)
	\shortans{$0,33$}
	\loigiai{
		\immini
		{
			Ta có $(\alpha)$ cắt các tia $Oy$, $Oz$ lần lượt lại các điểm $B$, $C$ không trùng $O$, suy ra $B(0;b;0)$, $C(0;0;c)$ với $b,c>0$.\\
			Phương trình mặt phẳng $(\alpha)\colon \dfrac{x}{1}+\dfrac{y}{b}+\dfrac{z}{c}=1$.\\
			Mặt phẳng $(P)$ có một véc-tơ pháp tuyến là $\overrightarrow{n}_{(P)}=(1;0;-1)$.\\
			Mặt phẳng $(\alpha)$ có một véc-tơ pháp tuyến là $\overrightarrow{n}_{(\alpha)}=\left(1;\dfrac{1}{b};\dfrac{1}{c}\right)$.\\
			Ta có $(P)\perp (\alpha)\Leftrightarrow \overrightarrow{n}_{(P)}\cdot\overrightarrow{n}_{(\alpha)}=0\Leftrightarrow 1-\dfrac{1}{c}=0\Leftrightarrow c=1.$\\
			Khi đó $(\alpha)\colon \dfrac{x}{1}+\dfrac{y}{b}+\dfrac{z}{1}=1\Leftrightarrow bx+y+bz-b=0$.\\
			Ta có
			\begin{eqnarray*}
				\mathrm{d}(O,(\alpha))=\dfrac{2}{3}&\Leftrightarrow&\dfrac{|-b|}{\sqrt{{2b^2}+1}}=\dfrac{2}{3} \\
				&\Leftrightarrow& 9b^2=4(2b^2+1)\\
				&\Leftrightarrow& \hoac{& b=2~\,\text{(nhận)}\\ & b=-2~\, \text{(loại).}}
			\end{eqnarray*}
		}
		{
			\begin{tikzpicture}[scale=0.7, font=\footnotesize, line join=round, line cap=round, >=stealth]
				\path 
				(0,0) coordinate (O)
				(1.2,-1.5) coordinate (A)
				(4,0) coordinate (B)
				;	
				\coordinate (C) at ($(O)+(0,3)$);
				\coordinate (x) at ($(O)!1.6!(A)$);
				\coordinate (y) at ($(O)!1.3!(B)$);
				\coordinate (z) at ($(O)!1.4!(C)$);
				\coordinate (I) at ($(A)!0.5!(B)$);
				\coordinate (J) at ($(I)!0.5!(C)$);
				\draw pic[draw=black,angle radius=0.16cm] {right angle = O--I--A}
				pic[draw=black,angle radius=0.16cm] {right angle = O--J--I};
				\draw[->](A) -> (x)node[below]{$x$};\draw[->](B) -> (y)node[below]{$y$}; \draw[->](C) -> (z)node[left]{$z$};
				\draw (C)--(O)--(A)--(B)--cycle (A)--(C)--(I);
				\draw[dashed] (B)--(O)--(I) (O)--(J);
				\foreach \p/\q in {A/190,B/-90,C/180,O/180}
				\fill[black] (\p) circle (1.0pt) ($(\p)+(\q:3mm)$) node{$\p$};
			\end{tikzpicture}
		}
		\noindent Suy ra $(\alpha)\colon 2x+y+2z-2=0$ và $B(0;2;0)$, $C(0;0;1)$.\\
		Ta có $V_{O.ABC}=\dfrac{1}{6}\cdot OA\cdot OB\cdot OC =\dfrac{1}{6}\cdot1\cdot 2\cdot 1= \dfrac{1}{3}$.
	}
\end{ex}
\begin{ex}%[2H5C3-4]
	Một nguồn âm phát ra sóng âm là sóng cầu. Khi gắn hệ trục tọa độ $O x y z$ (đơn vị trên mỗi trục là mét). Cường độ âm chuẩn tại điểm $I(3 ; 4 ; 5)$ là tâm của nguồn phát âm với bán kính $10 \mathrm{~m}$. Để kiểm tra một điểm ở vị trí $M(7 ; 10 ; 17)$ có nhận được cường độ âm phát ra tại $I$ hay không người ta sẽ tính khoảng cách giữa hai vị trí $I$ và $M$. Hỏi khoảng cách giữa hai vị trí $I$ và $M$ là bao nhiêu mét?
	\shortans{$14$}
	\loigiai{
	Ta có $IM=\sqrt{(7-3)^2+(10-4)^2+(17-5)^2}=14m$.
	}
\end{ex}
\Closesolutionfile{ans}
\indapan{6}{ans/ans-2-Otc5-De2-KQ}